% ============================================================
% Section 76: 六角密排结构因子与钠金属带间跃迁
% ============================================================
\section{固体理论：六角密排结构因子与钠金属带间跃迁}
\label{sec:hcp-structure-factor}

\begin{modelbox}[题目：HCP结构因子、导电性与BCC钠的带间跃迁阈波长]
\begin{enumerate}
    \item 六角密排结构的周期势 $V(\bm{r})=\sum_G V_G e^{i\bm{G}\cdot\bm{r}}$。证明对应于 $c$ 方向第一布里渊区边界面（法线沿 $\langle001\rangle$）的倒格矢 $G$，有 $V_G=0$。据此对经此面的价能隙能得出什么结论？
    \item 解释为什么在六角密排结构中，不能由单价原子组成绝缘体；
    \item 钠金属为体心立方结构，每个原子有一个价电子。求钠金属中带间跃迁（从最低能带到次高能带）阈波长的表达式，并估算其值（以 \AA 为单位）。已知费米波矢 $k_F$ 远小于布里渊区边界上离原点最近的波矢 $k_{ZB}$，钠的立方晶格边长 $a=4.23$ \AA。
\end{enumerate}
\end{modelbox}

\begin{tipbox}[考点快速分析]
\begin{itemize}
    \item \textbf{结构因子}：$S(\bm{G})=\sum_j e^{i\bm{G}\cdot\bm{d}_j}$，若 $S(G)=0$ 则 $V_G=0$
    \item \textbf{HCP基矢}：原胞含2个原子，$\bm{d}_A=0$，$\bm{d}_B=\frac{a}{2}\hat{\bm{i}}+\frac{\sqrt{3}a}{6}\hat{\bm{j}}+\frac{c}{2}\hat{\bm{k}}}$
    \item \textbf{能隙为零}：能带在布里渊区边界接触，系统呈金属性
    \item \textbf{带间跃迁}：$\hbar\omega_{\text{th}}=E(2k_{ZB}-k_F)-E(k_F)$
\end{itemize}
\end{tipbox}

\begin{derivationbox}[标准解题过程]
\textbf{Step 1: $c$ 方向布里渊区边界的 $V_G=0$}

六角密排（HCP）结构每个原胞含 2 个原子，坐标为
\begin{equation}
\bm{d}_A=0,\qquad
\bm{d}_B=\frac{a}{2}\hat{\bm{i}}+\frac{\sqrt{3}a}{6}\hat{\bm{j}}+\frac{c}{2}\hat{\bm{k}}.
\end{equation}

周期势的傅里叶系数 $V_G$ 正比于几何结构因子
\begin{equation}
S(\bm{G})=\sum_j e^{i\bm{G}\cdot\bm{d}_j}=1+e^{i\bm{G}\cdot\bm{d}_B}.
\end{equation}

对应于 $c$ 方向的第一布里渊区边界，法线沿 $\langle001\rangle$，相应的倒格矢为
\begin{equation}
\bm{G}=\frac{2\pi}{c}\,\hat{\bm{k}}.
\end{equation}

代入 $\bm{d}_B$：
\begin{equation}
\bm{G}\cdot\bm{d}_B=\frac{2\pi}{c}\cdot\frac{c}{2}=\pi
\quad\Longrightarrow\quad
S(\bm{G})=1+e^{i\pi}=1-1=0.
\end{equation}

因此 $V_G\propto S(\bm{G})=0$，即
\begin{equation}
\boxed{V_G=0.}
\end{equation}

\textbf{结论}：在近自由电子近似下，该布里渊区边界上的能隙
\begin{equation}
\Delta E=2|V_G|=0.
\end{equation}
即在 $c$ 轴方向，第一和第二能带在边界处简并，不发生分裂。

\textbf{Step 2: HCP 单价原子不能成绝缘体}

HCP 每个原胞含 2 个原子。若原子为单价，每个原胞贡献 2 个价电子。考虑自旋简并，每条能带可容纳 $2N$ 个电子（$N$ 为原胞数）。因此 $2N$ 个价电子恰好填满第一布里渊区对应的第一能带。

然而，由 (1) 的结果，沿 $c$ 轴方向（$\Gamma\to A$）第一布里渊区边界处能隙为零。这意味着第一能带和第二能带在该方向上\textbf{接触}，费米面穿过能带接触点。电子可无热激发地占据第二能带的底部，系统表现出\textbf{金属性}（或半金属性）。

\textbf{实例}：碱土金属 Be、Mg 为 HCP 结构，二价，本应填满第一能带为绝缘体，但因能带交叠，实际为金属。

\textbf{Step 3: BCC 钠的带间跃迁阈波长}

\textbf{物理图像}：带间跃迁是电子从最低能带（部分填充）吸收光子跃迁到次高能带。跃迁所需最小光子能量对应从费米能级到次高能带底部的能量差。

在近自由电子模型中，钠为 BCC 结构，其倒格子为 FCC。第一布里渊区边界上离原点最近的点为 $N$ 点（沿 $\langle110\rangle$ 方向），波矢大小
\begin{equation}
k_{ZB}=\frac{\sqrt{2}\,\pi}{a}.
\end{equation}

由于 $k_F\ll k_{ZB}$，费米球完全落在第一布里渊区内。带间跃迁阈能量对应于从 $k_F$ 跃迁到第二能带在布里渊区边界附近的状态。简化的判据为
\begin{equation}
\hbar\omega_{\text{th}}=E(2k_{ZB}-k_F)-E(k_F)
=\frac{\hbar^2}{2m}\bigl[(2k_{ZB}-k_F)^2-k_F^2\bigr]
=\frac{2\hbar^2}{m}k_{ZB}(k_{ZB}-k_F).
\end{equation}

\textbf{计算 $k_F$ 和 $k_{ZB}$}：
\begin{itemize}
    \item 钠为 BCC，每个惯用晶胞含 2 个原子，单价。电子密度 $n=2/a^3$。
    \item $k_F=(3\pi^2n)^{1/3}=(6\pi^2)^{1/3}/a\approx 3.898/a$
    \item $k_{ZB}=\sqrt{2}\,\pi/a\approx 4.442/a$
\end{itemize}

代入 $a=4.23$ \AA：
\begin{align}
k_F&\approx\frac{3.898}{4.23}\approx0.921\,\text{\AA}^{-1},\\
k_{ZB}&\approx\frac{4.442}{4.23}\approx1.050\,\text{\AA}^{-1},\\
k_{ZB}-k_F&\approx0.129\,\text{\AA}^{-1}.
\end{align}

阈波长 $\lambda_{\text{th}}=hc/(\hbar\omega_{\text{th}})$。利用 $\hbar c=1973\,\text{eV}\cdot\text{\AA}$，$mc^2=0.511\times10^6\,\text{eV}$：
\begin{align}
\lambda_{\text{th}}
&=\frac{\pi mc}{\hbar}\cdot\frac{1}{k_{ZB}(k_{ZB}-k_F)} \notag\\
&=\frac{mc^2}{\hbar c}\cdot\frac{\pi}{k_{ZB}(k_{ZB}-k_F)} \notag\\
&\approx\frac{0.511\times10^6}{1973}\cdot\frac{\pi}{1.050\times0.129}\notag\\
&\approx259\times23.3\approx\boxed{6030\,\text{\AA}.}
\end{align}
\end{derivationbox}

\begin{formulabox}[答案汇总]
\begin{center}
\begin{tabular}{ll}
\toprule
\textbf{问题} & \textbf{答案} \\
\midrule
(1) $V_G=0$ & 结构因子 $S(G)=1+e^{i\pi}=0$；能隙 $\Delta E=0$ \\
(2) 不能成绝缘体 & $c$ 轴方向能带接触，电子可自由跃迁，呈金属性 \\
(3) 阈波长 & $\lambda_{\text{th}}\approx 6030$ \AA \\
\bottomrule
\end{tabular}
\end{center}
\end{formulabox}

\begin{insightbox}[物理图像]
\begin{itemize}
    \item 晶体势的傅里叶分量由几何结构因子决定。特定晶面的``消光''（结构因子为零）导致该方向能隙消失，这是晶体对称性的直接后果。
    \item 能带接触使满带不成为绝缘体的充分条件。HCP 单价金属（如二价碱土金属 Be、Mg）因能带交叠而呈金属性。
    \item 钠的阈波长约 6030 \AA，处于可见光--近红外边界附近。这与碱金属在可见光区具有较高反射率（呈现金属光泽）的光学性质一致：光子能量低于带间跃迁阈值时，电子无法吸收光子跃迁到高能带，光子被反射而非吸收。
\end{itemize}
\end{insightbox}
